\section*{الفصل \ref{ch:b2:magmatic-series} --- المتتاليات الصهارية والأطر التكتونية}

\begin{solution}{exo:b2:magmatic-series:1}
(أ) $\ce{SiO2} = 49.60$، والقلويات $2.22 + 0.51 = 2.73$: بين 45 و52 وأدنى من 5،
\emph{بازلت} (B). (ب) 58.5 و$3.9 + 1.6 = 5.5$: بين 57 و63؛ والحدّ
الأعلى هناك يمتدّ من \babelsublr{(57, 5.9)} إلى \babelsublr{(63, 7.0)}، فيكون عند $5.9 + 1.1\times1.5/6 = 6.18$
من أجل 58.5؛ و5.5 تقع تحته: \emph{أنديزيت} (O2). (ج) 72.0 و$4.0 + 4.2 = 8.2$:
بعد الخط الرأسي عند 69 وفوق 8 وفوق الخط من \babelsublr{(69, 8)} إلى \babelsublr{(77, 0)}، الذي يكون
عند 5.0 من أجل 72: \emph{ريوليت} (R). (د) 47.0 و$4.5 + 2.5 = 7.0$: عند 47 يقع الحقل U1
بين الخط من \babelsublr{(45, 5)} إلى \babelsublr{(49.4, 7.3)}، عند $5 + 2.3\times2/4.4 = 6.05$،
والخط من \babelsublr{(45, 9.4)} إلى \babelsublr{(49.4, 7.3)}، عند $9.4 - 2.1\times2/4.4 = 8.45$. فالصخر
في U1: \emph{تيفريت} إذا كان في معياره أقل من \qty{10}{\percent} من الأوليفين،
و\emph{بازانيت} إذا كان أكثر.
\end{solution}

\begin{solution}{exo:b2:magmatic-series:2}
(أ) مجموع الأكاسيد 98.40 والمجموع مع المتطايرات 100.00؛
\ce{H2O} $1.40 < 2$ و\ce{CO2} $0.20 < 0.5$: فالتحليل مقبول. (ب) العامل
$100/98.40 = 1.01626$: \ce{SiO2} 52.44، \ce{TiO2} 1.22، \ce{Al2O3} 16.36،
\ce{FeO_T} 9.96، \ce{MgO} 6.71، \ce{CaO} 8.33، \ce{Na2O} 3.56، \ce{K2O} 1.42
(المجموع 100.00)؛ والقلويات $3.56 + 1.42 = 4.98$. (ج) قبل: 51.60 و4.90، تحت
حدّ البازلت--الأنديزيت البازلتي عند 52 مباشرة: \emph{بازلت}. بعد: 52.44 و4.98؛ وعند
52.44 يكون الحدّ الأعلى للحقل O1، من \babelsublr{(52, 5)} إلى \babelsublr{(57, 5.9)}، عند $5 + 0.9\times0.44/5
= 5.08$: \emph{أنديزيت بازلتي} (O1). فقد نقلت إعادة الحساب الصخر عبر
حدّ: والأسماء قرب خط تتوقّف عليها. (د) عند 52.44 يكون خط البوتاسيوم المنخفض عند $0.3 +
0.045\times4.44 = 0.50$ وخط البوتاسيوم العالي عند $1.2 + 0.095\times4.44 = 1.62$؛ و1.42
تقع بينهما: \emph{متوسط البوتاسيوم}.
\end{solution}

\begin{solution}{exo:b2:magmatic-series:3}
(أ) $Q + A + P = 90$: $Q = 25/90 = \qty{27.8}{\percent}$ و$100\,P/(A + P) =
100\times35/65 = 53.8$: الحقل 3b، \emph{مونزوغرانيت}، أي غرانيت. والبيوتيت
والهورنبلند لا يدخلان مجموع QAPF. (ب) $Q + A + P = 65$: $Q = 5/65 =
\qty{7.7}{\percent}$ (بين 5 و20) و$100\times55/60 = 91.7$ (فوق 90): الحقل
10*، ديوريت كوارتزي أو غابرو كوارتزي أو أنورثوزيت كوارتزي. \omterm{def:b2:magmatic-series:qapf}{ومخطّط QAPF} وحده لا
يختار بينها؛ بل يختار البلاجيوكليز والمعادن المافية (وهنا
\qty{35}{\percent} من الكلينوبيروكسين).
\end{solution}

\begin{solution}{exo:b2:magmatic-series:4}
خط البوتاسيوم المنخفض هو $0.3 + 0.045(\ce{SiO2} - 48)$، وخط البوتاسيوم العالي $1.2 +
0.095(\ce{SiO2} - 48)$. (أ) عند 50: 0.39 و1.39؛ و$0.40$ فوق الخط المنخفض بقليل:
\emph{متوسط البوتاسيوم}. (ب) عند 56: 0.66 و1.96: \emph{متوسط البوتاسيوم}. (ج) عند 62: 0.93 و2.53؛
$2.90 > 2.53$: \emph{عالي البوتاسيوم}. (د) عند 66: 1.11 و2.91؛ $1.00 < 1.11$: \emph{منخفض البوتاسيوم}.
\end{solution}

\begin{solution}{exo:b2:magmatic-series:5}
(أ) البازلت: $A + F + M = 3.2 + 14.5 + 4.5 = 22.2$؛ $A = 14.4$، $F = 65.3$، $M = 20.3$.
الأنديزيت: $5.05 + 6.1 + 3.8 = 14.95$؛ $A = 33.8$، $F = 40.8$، $M = 25.4$. (ب) الخط
لا يتجاوز أبدًا $F = 54.5$؛ فالبازلت، الذي فيه $F = 65.3$، يقع في جهة F
أيًّا كانت قيمة $A$: \emph{ثولييتي}، بازلت حديدي من اتجاه فينر. أمّا
الأنديزيت، فعند $A = 33.8$ يقع الخط بين \babelsublr{(49.6, 35.0)} و\babelsublr{(52.5, 30.0)}، عند $F =
49.6 + 2.9\times1.2/5 = 50.3$؛ $40.8 < 50.3$: \emph{كلسي قلوي}. (ج) الماغنتيت،
الذي يحمل الحديد الحديديك خارج صهارة مؤكسدة منذ البداية ويمنع $F/M$ من
الارتفاع بالسرعة التي يرتفع بها في الثولييت.
\end{solution}

\begin{solution}{exo:b2:magmatic-series:6}
بازلت MORB: $\mathrm{Nb_N} = 4/0.658 = 6.08$، $\mathrm{La_N} = 4.12/0.648 = 6.36$، والنسبة
0.956. \omterm{def:b2:elements-in-the-earth:morb-oib}{بازلت الجزر المحيطية}: $38/0.658 = 57.8$، $30.8/0.648 = 47.5$، والنسبة 1.215. بازلت
الأقواس: $3/0.658 = 4.56$، $6.1/0.648 = 9.41$، والنسبة 0.484. بازلت الأقواس وحده فيه
شذوذ نيوبيوم، أي حوض عند Nb في المخطّط العنكبوتي: فالنيوبيوم مستنزَف مقارنةً
بجاره La: لقد تلقّت صهارة القوس عناصر من اللوح، لكن قليلًا من
النيوبيوم. وبازلت MORB شبه مستوٍ؛ \omterm{def:b2:elements-in-the-earth:morb-oib}{وبازلت الجزر المحيطية} أغنى
بالنيوبيوم منه باللانثانوم.
\end{solution}

\begin{solution}{exo:b2:magmatic-series:7}
(أ) $\mathrm{Ti} = 1.6\times0.5993\times10^4 = \qty{9589}{ppm}$؛ \babelsublr{Ti/100 = 95.9}، و\babelsublr{Zr
110}، و$3\mathrm Y = 96$؛ والمجموع 301.9؛ والإحداثيات 31.8 و36.4 و31.8: الحقل B. (ب) Ti
16\,780، ومنه 167.8 و220 و87؛ والمجموع 474.8؛ 35.3 و46.3 و18.3: الحقل D. (ج) Ti 4794، ومنه
47.9 و50 و60؛ والمجموع 157.9؛ 30.3 و31.7 و38.0: الحقل B. فالبازلت (ب) بازلت داخل الصفائح
(لجزيرة محيطية أو قارة)؛ و(أ) و(ج)، في الحقل B، قد يكونان بازلت MORB أو ثولييتات أقواس
أو بازلتات كلسية قلوية: فالمخطّط لا يستطيع الحسم، ويلزم مخطّط آخر (مخطّط
عنكبوتي، أو شذوذ Nb).
\end{solution}

\begin{solution}{exo:b2:magmatic-series:8}
(أ) بازلت MORB $142/164 = 0.866$، وبازلتات الجزر المحيطية $4/173 = 0.023$، وبازلتات الأقواس $101/201 =
0.502$. (ب) مع احتمالات مسبقة متساوية، تعطي قاعدة بايز $0.866/(0.866 + 0.023 + 0.502) =
0.866/1.391 = 0.623$. (ج) أن البازلت يُستبعد جدًا أن يكون بازلت جزر محيطية
(الاحتمال $0.023/1.391 = 0.017$)، لكنه لا يقول هل هو بازلت MORB أم بازلت أقواس:
نحو ثلاث فرص من خمس مقابل اثنتين.
\end{solution}

\begin{solution}{exo:b2:magmatic-series:9}
المولات لكل \qty{100}{g}: \ce{Na2O} $5.5/61.98 = 0.08874$، \ce{K2O} $4.5/94.20 =
0.04777$، \ce{Al2O3} $10.5/101.96 = 0.10298$. (أ) $(0.08874 + 0.04777)/0.10298 =
0.13651/0.10298 = 1.326 > 1$: فالريوليت \emph{فوق قلوي}. (ب) يستطيع الفلسبار
أن يأخذ $0.10298$ من مولات الأكسيد القلوي البالغة $0.13651$؛ والفائض، $0.03353$، هو
$0.03353/0.13651 = 0.246$، أي ربع القلويات، التي يجب أن تدخل معادن أخرى.
\end{solution}

\begin{solution}{exo:b2:magmatic-series:10}
(أ) $s = (\lambda_F - \lambda_M)/(\lambda_F - \lambda_A)$: من أجل $\lambda_F = 0.5$، $(0.5
- 0.8)/(0.5 - 0.4) = -3.0$؛ ومن أجل $0.7$، $(-0.1)/(0.3) = -0.333$. (ب) حلّ
$s(\lambda_F - 0.4) = \lambda_F - 0.8$ يعطي $\lambda_F = (0.4s - 0.8)/(s - 1)$:
آيسلندا، $s = -1.10$: $(-0.44 - 0.8)/(-2.10) = 0.590$؛ الكاسكيد، $s = -0.72$:
$(-0.288 - 0.8)/(-1.72) = 0.633$. (ج) صهارة الكاسكيد، ذات القيمة الأكبر للمقدار
$\lambda_F$، تفقد حديدها أسرع: إذ يُخرجه الماغنتيت. (د) ينمو $\ln(F/M)$
بالمعدل $\lambda_M - \lambda_F$ و$\ln(A/F)$ بالمعدل $\lambda_F - \lambda_A$؛ فلكل
وحدة يكسبها $\ln(A/F)$ يكسب $\ln(F/M)$ المقدار $-s$: 1.10 في آيسلندا و0.72 في
الكاسكيد. فترتفع $F/M$ على امتداد الاتجاهين، لكن أسرع في الاتجاه الثولييتي، أي اتجاه
فينر.
\end{solution}

\begin{solution}{exo:b2:magmatic-series:11}
(أ) BHVO-2، المولات لكل \qty{100}{g}: \ce{SiO2} 0.82557، \ce{Al2O3} 0.13182، \ce{FeO}
0.15521، \ce{MgO} 0.18007، \ce{CaO} 0.20328، \ce{Na2O} 0.03580، \ce{K2O} 0.00545.
الأورثوكليز 0.00545 والألبيت 0.03580؛ و\ce{Al2O3} الباقي $0.13182 - 0.04125 = 0.09057$،
كله في الأنورثيت مع القدر نفسه من \ce{CaO}؛ و\ce{CaO} الباقي $0.11271$ في الديوبسيد؛
و$(\ce{Mg,Fe})\ce{O}$ الباقي $0.33528 - 0.11271 = 0.22257$ في الهيبرستين. السيليكا
اللازمة: $6\times0.00545 + 6\times0.03580 + 2\times0.09057 + 2\times0.11271 + 0.22257 =
0.87663$؛ والمتاح 0.82557؛ فالنقص 0.05106. وكل $(\ce{Mg,Fe})\ce{O}$ يتحوّل من
الهيبرستين إلى الأوليفين يوفّر نصف سيليكا، وبحدّ أقصى $0.22257/2 = 0.11129$: فيكفي تحويل
$0.10212$ منه، ويبقى $0.12045$ في الهيبرستين. ففي المعيار أوليفين
وهيبرستين، ولا كوارتز ولا نيفلين: فالصخر مشبّع بالسيليكا بمعنى
التعريف، \emph{ثولييت أوليفيني} في معياره الهيبرستين، من \omterm{def:b2:magmatic-series:series}{المتتالية
الثولييتية}، كما تقذف مرحلة الدرع في بركان هاوايي. (ب) المولات: 0.75732، 0.14712، 0.16008، 0.21092، 0.17832، 0.06131،
0.01592. الأورثوكليز 0.01592، والألبيت 0.06131، والأنورثيت 0.06988، والديوبسيد 0.10843، والهيبرستين $0.37100 - 0.10843 =
0.26257$؛ والسيليكا اللازمة $0.09552 + 0.36786 + 0.13976 + 0.21686 + 0.26257 = 1.08257$؛
والمتاح 0.75732؛ فالنقص 0.32525. وتحويل الهيبرستين كله إلى أوليفين يوفّر
0.13129، فيبقى 0.19396؛ وكل مول من \ce{Na2O} يتحوّل من الألبيت إلى النيفلين
يحرّر 4 مولات من \ce{SiO2}، ومنه $0.19396/4 = 0.04849$
من مولات \ce{Na2O} الألبيت البالغة 0.06131 تصير نيفلينًا. ففي المعيار نيفلين:
والصخر دون مشبّع بالسيليكا، \emph{بازلت قلوي}، من \omterm{def:b2:magmatic-series:series}{المتتالية القلوية}.
\end{solution}

\begin{solution}{exo:b2:magmatic-series:12}
(أ) بازلتات الأقواس $116/141 = 0.823$، وبازلت MORB $21/101 = 0.208$، وبازلتات الجزر المحيطية
$3/137 = 0.022$. (ب) $0.823/(0.823 + 0.208 + 0.022) = 0.823/1.053 = 0.782$.
(ج) البسط $0.1\times0.823 = 0.0823$؛ والمقام $0.0823 + 0.6\times0.208 +
0.3\times0.022 = 0.0823 + 0.1248 + 0.0066 = 0.2137$؛ والاحتمال $0.385$. (د) الاختبار
نفسه على الصخر نفسه يعطي 0.78 أو 0.39 بحسب مدى شيوع كل إطار
حيث وُجد الصخر: \omterm{def:b2:magmatic-series:discrimination}{فمخطّط التمييز} يعطي أرجحية، ويجب أن يوفّر
السياق الجيولوجي الباقي.
\end{solution}

\begin{solution}{pb:b2:magmatic-series:1}
\noindent\textbf{الجزء الأول.}

\textbf{1.} الأكاسيد: الحمم 1 99.17، الحمم 2 98.65، الحمم 3 98.20؛ ومع المتطايرات
99.97 و99.80 و99.45، وكلها قريبة من 100. والماء (0.70، 1.10، 1.20) أدنى من
\qty{2}{\percent} وثاني أكسيد الكربون (0.10، 0.05، 0.05) أدنى من \qty{0.5}{\percent}: فالثلاثة
كلها مقبولة.

\textbf{2.} العامل $100/99.17 = 1.00837$: \ce{SiO2} 50.82، \ce{Na2O} 2.92، \ce{K2O}
0.63، والقلويات 3.55.

\textbf{3.} الحمم 2: $100/98.65 = 1.01368$، \ce{SiO2} 59.40، \ce{Na2O} 3.95، \ce{K2O}
1.17، والقلويات 5.12. الحمم 3: $100/98.20 = 1.01833$، \ce{SiO2} 66.50، \ce{Na2O} 4.48،
\ce{K2O} 1.88، والقلويات 6.36.

\textbf{4.} الحمم 1 (50.82 و3.55): بازلت. الحمم 2 (59.40 و5.12؛ والحدّ الأعلى للحقل O2
هناك عند $5.9 + 1.1\times2.40/6 = 6.34$): أنديزيت. الحمم 3 (66.50 و6.36؛ بين 63
والخط من \babelsublr{(69, 8)} إلى \babelsublr{(77, 0)}، وتحت الخط من \babelsublr{(63, 7)} إلى \babelsublr{(69, 8)}): داسيت.

\textbf{5.} القيم الخام: (50.4 و3.52)، (58.6 و5.05)، (65.3 و6.25): الحقول نفسها.
وإعادة الحساب تنقل كل نقطة بما لا يزيد على \qty{1.2}{\percent} من السيليكا ولا تغيّر
شيئًا هنا، بخلاف \cref{exo:b2:magmatic-series:2}.

\textbf{6.} \omterm{def:b2:magmatic-series:qapf}{مخطّط QAPF}، من \omterm{def:b2:magmatic-series:qapf}{التركيب المعدني النمطي} المقيس. والصخر الجوفي ذو التركيب
الأنديزيتي غني بالبلاجيوكليز، وفيه قليل من فلسبار القلويات أو الكوارتز: فيُتوقّع أن
يقع في الحقل 10 أو 10*، ديوريت أو ديوريت كوارتزي (وتركيبه النمطي يحسم).

\noindent\textbf{الجزء الثاني.}

\textbf{7.} الحمم 1: الخطان عند 0.43 و1.47، و$\ce{K2O} = 0.63$: متوسطة البوتاسيوم. الحمم 2: 0.81
و2.28، و1.17: متوسطة البوتاسيوم. الحمم 3: 1.13 و2.96، و1.88: متوسطة البوتاسيوم.

\textbf{8.} \ce{FeO_T} و\ce{MgO} بعد استبعاد المتطايرات: 8.67 و7.97؛ 6.18 و3.85؛ 4.07
و1.63. الحمم 1: $3.55 + 8.67 + 7.97 = 20.19$؛ $A = 17.6$، $F = 42.9$، $M = 39.5$.
الحمم 2: 15.15؛ 33.8 و40.8 و25.4. الحمم 3: 12.06؛ 52.7 و33.7 و13.5.

\textbf{9.} الحمم 1، $A = 17.6$: الخط بين \babelsublr{(42.8, 17.1)} و\babelsublr{(46.9, 18.1)} يقع عند
$F = 42.8 + 4.1\times0.5/1.0 = 44.9$؛ $42.9 < 44.9$، الجهة الكلسية القلوية، قرب
الخط. الحمم 2، $A = 33.8$: $F = 49.6 + 2.9\times1.2/5 = 50.3$؛ $40.8 < 50.3$،
كلسية قلوية. الحمم 3، $A = 52.7$: $F = 40.8 - 2.9\times2.7/5 = 39.2$؛ $33.7 < 39.2$،
كلسية قلوية.

\textbf{10.} $\ce{FeO_T}/\ce{MgO}$: $8.67/7.97 = 1.09$، $6.18/3.85 = 1.61$، $4.07/1.63 =
2.50$. ترتفع النسبة، كما في كل بازلت يتمايز، لكن الحديد نفسه
ينخفض من 8.67 إلى 4.07، و$F$ ينخفض على المثلث من 42.9 إلى 33.7: فلا إثراء
بالحديد. والمجموعة الثولييتية كانت ستصعد أولًا نحو F، كما يفعل البازلت الحديدي في
\cref{exo:b2:magmatic-series:5}.

\textbf{11.} \omterm{def:b2:magmatic-series:series}{المتتالية الكلسية القلوية}. فالماغنتيت، الذي يتبلور من صهارة مؤكسدة،
أزال الحديد مع المغنيسيوم؛ والأنواع الكيميائية المؤكسدة تأتي من اللوح الذي يفقد ماءه في
نطاق اندساس.

\noindent\textbf{الجزء الثالث.}

\textbf{12.} $\mathrm{Nb_N} = 2.0/0.658 = 3.04$، $\mathrm{La_N} = 7/0.648 = 10.80$،
$\mathrm{Th_N} = 1.1/0.0795 = 13.84$.

\textbf{13.} $3.04/10.80 = 0.281$، أدنى بكثير من القيمتين الوسيطتين لبازلت MORB (0.956)
وبازلتات الجزر المحيطية (1.215)، بل أدنى من القيمة الوسيطة لبازلت الأقواس (0.484): أي حوض
نيوبيوم عميق.

\textbf{14.} $13.84/3.04 = 4.55$؛ ولبازلت MORB الوسيط في المخطّط العنكبوتي $2.64/6.08 =
0.43$. فالثوريوم، كغيره من العناصر المُثراة في بازلتات الأقواس، أضافه
اللوح المندسّ.

\textbf{15.} \ce{TiO2} بعد استبعاد المتطايرات $0.95\times1.00837 = 0.958$؛ $\mathrm{Ti} =
0.958\times5993 = \qty{5741}{ppm}$؛ \babelsublr{Ti/100 = 57.4}، و\babelsublr{Zr 70}، و$3\mathrm Y = 54$، والمجموع 181.4:
31.6 و38.6 و29.8، في الحقل B، الذي يتقاسمه بازلت MORB وثولييتات الأقواس والبازلتات الكلسية القلوية.
وهو يستبعد بازلتًا داخل الصفائح لكنه لا يحسم بين ظهرة وقوس.
والمخطّط بُني للبازلتات: ففي الأنديزيت والداسيت أزالت الأكاسيدُ التيتانيومَ
وركّز التبلورُ الزركونيوم، وحقولهما (الحمم 2
تقع في C، والحمم 3 خارج كل الحقول) لا معنى لها.

\textbf{16.} متتالية كلسية قلوية متوسطة البوتاسيوم، مع حوض نيوبيوم عميق ونسبة
$\mathrm{Th_N}/\mathrm{Nb_N}$ مرتفعة، على بركان قذف من البازلت إلى الداسيت: فالبركان
نما في قوس بركاني فوق نطاق اندساس.

\noindent\textbf{الجزء الرابع.}

\textbf{17.} $C_L = C_0/F$ لعنصر غير متوافق تمامًا: $F = 0.63/1.17 =
0.538$.

\textbf{18.} الروبيديوم $10/20 = 0.500$؛ والثوريوم $1.1/2.1 = 0.524$. (العناصر
النزرة، المعطاة كما حُلّلت، كانت ستتغيّر بأقل من \qty{1}{\percent} بإعادة
الحساب.)

\textbf{19.} $C_L/C_0 = F^{D - 1}$: $125/70 = 1.786 = F^{-0.9}$، $F = 1.786^{-1/0.9} =
0.525$.

\textbf{20.} من البوتاس، $F = 1.17/1.88 = 0.622$ من الحمم 2 إلى الحمم 3. السترونشيوم:
$380/480 = 0.792 = 0.622^{D - 1}$، ومنه $D - 1 = \ln 0.792/\ln 0.622 = -0.234/-0.474 =
0.49$ و$D \approx 1.49$: فقد صار السترونشيوم متوافقًا. فالبلاجيوكليز، الذي يأخذ
السترونشيوم، كان يتبلور بوفرة.

\textbf{21.} من البوتاس، $F = 0.63/1.88 = 0.335$ (الروبيديوم $10/32 = 0.31$، والثوريوم
$1.1/3.4 = 0.32$): فنحو ثلثي البازلت، $1 - 0.335 = 0.665$، تبلور
قبل الداسيت.

\textbf{22.} البوتاس المُستوعَب كان سيرفع \ce{K2O} في الحمم 2 فوق ما
يعطيه التبلور وحده، فتخرج $F = 0.63/1.17$ منخفضة أكثر من اللازم، وكتلة
البلّورات مرتفعة أكثر من اللازم.

\textbf{23.} $(0.538 + 0.500 + 0.524 + 0.525)/4 = 2.087/4 = 0.522$؛ والبلّورات $1 - 0.522
= 0.478$ من البازلت، أي \textbf{\qty[text-series-to-math=true]{48}{kg} من
البلّورات لكل \qty{100}{kg} من البازلت} لصنع الأنديزيت.
\end{solution}
